\documentclass[12pt,letterpaper]{article}

%Packages
\usepackage{amsmath}
\usepackage{amsthm}
\usepackage{amsfonts}
\usepackage{amssymb}
\usepackage{amscd}
\usepackage{dsfont}
\usepackage{enumerate}
\usepackage{fancyhdr}
\usepackage{mathrsfs}
\usepackage{bbm}
\usepackage{framed}
\usepackage{mdframed}
\usepackage{cancel}
\usepackage{float}
\usepackage{mathtools}
%\usepackage[]{mcode}
\usepackage{graphicx}
\usepackage{tikz}
\usetikzlibrary{automata,arrows}
\usepackage{verbatim}

%Page formatting
\usepackage[letterpaper,voffset=-.5in,bmargin=3cm,footskip=1cm]{geometry}
\setlength{\parindent}{0.0in}
\setlength{\parskip}{0.1in}
\allowdisplaybreaks
\headheight 15pt
\headsep 10pt

% Common Commands
\newcommand\N{\mathbb N}
\newcommand\Z{\mathbb Z}
\newcommand\R{\mathbb R}
\newcommand\Q{\mathbb Q}
\newcommand\lcm{\operatorname{lcm}}
\newcommand\setbuilder[2]{\ensuremath{\left\{#1\;\middle|\;#2\right\}}}
\newcommand\E{\operatorname{E}}
\newcommand\V{\operatorname{V}}
\newcommand\Pow{\ensuremath{\operatorname{\mathcal{P}}}}

\DeclarePairedDelimiter\ceil{\lceil}{\rceil}
\DeclarePairedDelimiter\floor{\lfloor}{\rfloor}

% 22 style
\newcommand\hint[1]{\textbf{Hint}: #1}
\newcommand\note[1]{\textbf{Note}: #1}
\newenvironment{22enumerate}{\begin{enumerate}[a.]\itemsep0em}{\end{enumerate}}
\newenvironment{22itemize}{\begin{itemize}\itemsep0em}{\end{itemize}}
\fancypagestyle{firstpagestyle} {
  \renewcommand{\headrulewidth}{0pt}%
  \lhead{\textbf{CSCI 0220}}%
  \chead{\textbf{Discrete Structures and Probability}}%
  \rhead{Littman}%
}

\pagestyle{fancyplain}

\newcommand\EX{\mathds{E}}
\newcommand\PR{\mathds{P}}
\newcommand\zlam{Z_\lambda}
\DeclareMathOperator*{\argmin}{arg\,min}

%%%%%%%%%%%%%%%% COMPILE WITH \soltrue or \solfalse %%%%%%%%%%%%%%%%%%%%%%%%%%%%%
\newif\ifsol
\soltrue  % SOLUTIONS
%\solfalse % NO SOLUTIONS
%%%%%%%%%%%%%%%%%%%%%%%%%%%%%%%%%%%%%%%%%%%%%%%%%%%%%%%%%%%%%%

%%%%%%%%%%%%%%%% solution help %%%%%%%%%%%%%%%%%%%%%
\newcommand{\solu}[2]{ \begin{mdframed} \ifsol #2 \else \vspace{#1} \fi \end{mdframed} }
\newcommand{\solt}{ \ifsol (T) \fi }
\newcommand{\solf}{ \ifsol (F) \fi }
\newcommand{\solm}[1]{\ifsol \textit{(#1)} \fi}


\begin{document}
  \thispagestyle{firstpagestyle}
  \begin{center}
    {\large \textbf{Recitation 8}}
    
    {\large Counting and Intro to Probability}

  \end{center}
  
	\section*{Counting Arguments}
	
	A counting argument shows that the LHS (lefthand side) and the RHS (righthand side) of some equation count the same thing. Instead of using algebraic manipulation, we explain why both sides enumerate the elements of some set, just in different ways.
	
% 	That is, we think about a set, and we come up with a story that corresponds to the LHS that describes an enumeration of the elements of that set, and then we come up with a different story that corresponds to the RHS that describes an enumeration of the elements of that very same set. 
	
	For instance, consider the following identity.
    $$\binom{n}{k} = \binom{n-1}{k-1} + \binom{n-1}{k}$$
    
    Let $S$ be a set with $n$ elements.
    
    \begin{itemize}
        \item The LHS counts the number of ways to form a subset of $S$ size $k$.
        \item The RHS also counts the number of subsets of size $k$. It partitions it into two parts:
        \begin{itemize}
            \item $k$-sized subsets \textbf{with} a fixed element $x$. Since we already know $x$ is in the subset, so we just want to pick $k-1$ more elements from the $n-1$ remaining elements.
            \item $k$-sized subsets \textbf{without} $x$. We know we can't pick $x$ for our subsets, so we just choose $k$ elements from $n-1$ other elements in $S$.
        \end{itemize}
    \end{itemize}
    
	
	\textbf{Task 1}
	
	Use counting arguments to prove the following identities:
	\begin{22enumerate}
	    \item $$\binom{n}{k} = \binom{n}{n - k}$$
	    \solu{4cm}{The LHS is the number of ways to choose $k$ elements from a set of size $n$, which can be equivalently seen as the number of ways to choose $n-k$ elements \emph{not} to include, which is what the RHS counts.}
	   
    	\item $$\sum_{k = 0}^{n} \binom{n}{k} = 2^n$$
        \solu{4cm}{The LHS is the number of subsets of each size of a set of size $n$ added together. In other words, we can count the number of subsets with 0 elements ($n \choose 0$), the number of subsets with 1 element ($n \choose 1)$), and so on. If we add them up, we will get $\sum_{k = 0}^{n} \binom{n}{k}$.
        
        Alternatively, the RHS is the total number of subsets of a size of size $n$. This is because, for each element, we can choose to either contain it in the subset, or not contain it in the subset. Thus, for each of the $n$ elements, we have $2$ options: include or don't include. That gives us $2 \cdot 2 \cdot ... \cdot 2 = 2^n$ different subsets.
        
        Thus, $\sum_{k = 0}^{n} \binom{n}{k} = 2^n$.}
        
        \item \textit{Optional: }
        
        $$\binom{n}{m}\binom{m}{k} = \binom{n}{k}\binom{n - k}{m - k}$$
        
        \hint{Try this with small numbers where $k < m < n$.}
        \solu{4cm}{The LHS is the number of ways to choose $m$ members of a general committee, and then choose $k$ special members from that $m$ for a subcommittee. The RHS is the number of ways to choose $k$ members of a special subcommittee, and then the number of ways to include the unpicked members in the general committee of size $m$, but it only has $m - k$ spots left.}
    \end{22enumerate}

\textbf{Task 2}

Below is a table of counting problems involving putting $k$ balls in $n$ distinct bins, and seven expressions. Match each problem with its solution.
\vspace{0.5cm}
\begin{figure}[h]
    \centering
    \begin{tabular}{p{2.5cm}|c|c|c}
    & No restrictions & At \textbf{most} 1 ball per bin & At \textbf{least} 1 ball per bin  \\
    & & \small assume $k \leq n$ & \small assume $k \geq n$\\
    \hline
    Identical balls & \small \textcircled{1}& \small \textcircled{3} & \small \textcircled{5} \textit{Optional} \\
    &  &\\
    \hline
    Distinct balls &\small \textcircled{2} &\small \textcircled{4} & \small \textcircled{6} \textit{Optional} \\
    & &\\
\end{tabular}
\begin{tabular}{ccccccc}
    \\
    \textbf{A} & \textbf{B} & \textbf{C} & \textbf{D} & \textbf{E} & \textbf{F} & \textbf{G}\\
    \large $\binom{n}{k}$ & \large $n^k$ & \large $\binom{k-1}{n-1}$ & \large $\frac{n!}{(n-k)!}$ & \large $\binom{k+n-1}{k}$ & \large $n^k-\sum_{i=1}^{n-1} \binom{n}{i} (n-i)^k (-1)^{i+1}$& $n!{{k}\choose{n}}n^{k-n}$ \\
\end{tabular}
\end{figure}
\solu{5cm}{
\begin{enumerate}[1.]
    \item This is stars and bars, so it's E.
    \item Each of the $k$ balls have $n$ options, so it's B.
    \item We have $n$ bins and $k$ of them will have 1 ball in them. There is a bijection between this situation and $n$-bit binary strings with $k$ ones and $n-k$ zeros. So, this is choosing which digits will have 1s, which is A.
    \item For each of the options from question 3, which is $\binom{n}{k}$, there are $k!$ ways of scrambling the balls in sequence. So we have $\frac{n!}{k!(n-k)!}k!$, which cancels out to D. Note this is also the permutation $P(n, k)$.
    \item This can be reduced to stars and bars, but we pre-distribute $n$ balls to the $n$ bins to make sure there's 1 in each bin. Then, with the $k-n$ balls left, we distribute it into $k$ bins, which is $\binom{k-n+n-1}{n-1}=\binom{k-1}{n-1}$, which is C.
    \item  We'll count the ways there is a bin with \textit{no} balls, using inclusion-exclusion, and find the complement of that. Let $E_j$ be the event that the $j$-th bin is empty, and we want to find $|E_1 \cup \cdots \cup E_n|$. Each $E_j$ is of size $(n-1)^k$, since the $k$ balls can go to any of the $n-1$ other bins. Then we take out the pairwise intersections, each of which are of size $(n-2)^k$, and so on to $(n-i)^k$ generally. The number of $i$-way intersections is $\binom{n}{i}$, which we also have to multiply by. To include all of this in a summation, along with using $(-1)^{i+1}$ to simulate the alternating addition/subtraction. Note that the last term, $i=n$, is always zero (why?). This is expression F.
    
    \textit{Note:} This is equivalent to asking how many surjective functions there are from a domain of size $k$ to a codomain of size $n$.
\end{enumerate}
}
\subsection*{Checkpoint 1 - Call over a TA}
\newpage

\subsection*{Inclusion/Exclusion}

Say we have two sets, $A$ and $B$, and we want to know how many elements there are in their union, $A \cup B$. If $A$ and $B$ have no elements in common, we can compute this just by adding the cardinality of each, $|A| + |B|$. However, this approach will not work in general.

For instance, if $A = \{1, 2\}$ and $B = \{1, 3\}$, then $A \cup B = \{1, 2, 3\}$. $|A \cup B| = 3$, but $|A| + |B| = 2 + 2 = 4$. The problem is that we've double counted the elements that are in both $A$ and $B$, that is, element 1. To fix this problem, we should subtract $|A \cap B|$.

\begin{center} \includegraphics[scale=.5]{venn.png} \end{center}

The resulting Inclusion/Exclusion property for two sets is: $$|A \cup B| = |A| + |B| - |A \cap B|.$$

Let's try generalizing this idea to three sets, $A$, $B$, and $C$!
% if we first add the cardinality of each set, . 

\begin{center} \includegraphics[scale=.5]{venn3.png} \end{center}

If we tried adding $|A|+|B|+|C|$, which regions would we double count? Which regions would we triple count?
\solu{1cm}{$w, x, y$ get double counted and $z$ gets triple counted.}

As it turns out, the final formula for $|A \cup B \cup C|$ is the following. Convince yourself that it works!
$$|A \cup B \cup C| = |A|+|B|+|C| - |A \cap B| - |B \cap C| - |A \cap C| + |A \cap B \cap C|.$$

Going further, we can repeat this process for any number of sets, alternating between adding and subtracting the sizes of sets.

\textbf{Task 3}
\begin{22enumerate}
    \item Let $S=\{1, 2, 3, 4, 5\}$.
    \begin{enumerate}[i.]
        \item How many permutations of $S$ contain the sequence 24?
        \solu{2cm}{Treat 24 as a single units, so we're left with permuting \{1, 3, 5, 24\}. There are 4! permutations.}
        \item How many permutations of $S$ contain the sequence 52?
        \solu{1cm}{Same as above, 4!.}
        \item How many permutations of $S$ contain the sequence 24 \textit{or} 52?
        \solu{2cm}{Let $A$ be the set of permutations where $24$ appears and $B$ be the set of permutations where $52$ appears. $A \cap B$ will be the set of permutations where $524$ appears. If we treat this as a unit, we permute 3 elements, $|A \cap B| = 3!$. Thus, $|A \cup B| = |A| + |B| - |A \cap B| = 4! + 4! - 3!=42$.}
    \end{enumerate}
    \item Let $X$ and $Y$ be sets such that $|X| = 8$ and $Y = \{a, b, c\}$.
    \begin{enumerate}[i.]
        \item How many functions from $X \rightarrow Y$ do not map to $a$?
        \solu{1cm}{Each of the 8 elements in $X$ have two options of where to map to: $b$ or $c$. So, it's $2^8=256$.}
        \item How many functions from $X \rightarrow Y$ do not map to $a$ and also do not map to $b$? 
        \solu{1cm}{There's only 1 option, so it's $1^8=1$.}
        \item How many functions from $X \rightarrow Y$ are \textit{not} surjective?
        
        \hint{A function is not surjective if nothing maps to $a$, nothing maps to $b$, or nothing maps to $c$.}
        \solu{4cm}{Let the set of functions not mapping to $a$ be $A$, not mapping to $b$ be $B$, and not mapping to $c$ be $C$. \\
        $|A|=|B|=|C|=2^8$, as established above, since the argument for $i$ can be made to all three sets. \\
        $|A \cap B| = |B \cap C| = |A \cap C| = 1$, as the argument for ii can also be generalized. \\
        $A \cap B \cap C$ is functions from $X$ to $\{\}$, which is no functions. \\
        The total number of non-surjective functions is $|A \cup B \cup C| = 3 * 2^8 - 3 * 1 + 0=765$.} 
    \end{enumerate}
    
    
    
\end{22enumerate}

\subsection*{The Pigeonhole Principle}

Let's say we have $n$ pigeons who are trying to fit in $n-1$ holes.

 \includegraphics[scale=0.25]{pigeons.jpg}

It isn't possible for each pigeon to get its own hole: at least two of them are going to have to share. It could be the case that they are all in the same hole, or, like the picture above, all but 2 pigeons get their own hole, or anything in between.

This is the Pigeonhole Principle: in general, if we are assigning $n$ objects to $m$ categories, where $n > m$, there is at least one category that has more than one object assigned to it.

Solve the following problems using the Pigeonhole Principle:

\textbf{Task 4}
\begin{22enumerate}
  \item Celeste is pulling socks out of her drawer. She only has four types of socks: solid, striped, polka-dotted, and ones with planets on them. What is the minimum number of socks Celeste should pull out to ensure she has a pair? 
  
  \solu{1 cm}{We are asking what the size of the domain of a function needs to be, where the function is from socks to types, such that there is some style with at least 2 socks that map to it. The answer is therefore 5.}
  
  \item  There are $n > 2$ astronauts having a party on Mars. Throughout the night, they dance with each other in pairs.

		\begin{enumerate}[i.]
			\item The minimum number of total dance partners someone can have is 0 (they didn't dance with anyone). What is the maximum number of dance partners one can have?
			
			\solu{2cm}{$n - 1$, dancing with everyone else}

			\item Prove that at least 2 astronauts have the same number of dance partners by the end of the night. (Come back to this question later if you get stuck.)
			
			\solu{2cm}{The possible number of dance partners are $[0, n - 1]$. Trying to apply the pigeonhole principle directly doesn't work, since it seems like all $n$ astronauts could have all $n$ possible numbers of partners. But, if someone has $0$ dance partners, then no one can have $n - 1$. So, the choices are either $[1, n-1]$ or $[0, n-2]$, each of which results in $n - 1$ possible options for the number of dance partners. Since we have $n$ astronauts, two people must have the same number of partners.}

		\end{enumerate}
 \begin{comment}
 \item Suppose $S$ is a set of $n + 1$ integers. Prove that there exist distinct $a, b \in S$ such that $a - b$ is a multiple of $n$.
 
 \solu{3 cm}{Consider the differences $a_{n + 1} - a_1$ .... $a_{n + 1} - a_n$. If one of these differences has a remainder of $0$ mod $n$, we are done. Otherwise, two of them have the same remainder:  $a_{n + 1} - a_i$ and  $a_{n + 1} - a_j$ where $i \neq j$. Subtract them from each other, and we are done.}
\end{comment}
 \item \textit{Optional:} Given any 5 points inside a square with side length 2, there is always a pair whose distance apart is at most square root of 2. 
 
 \solu{6 cm}{Partition the square into 4 unit squares. Then, there is a square with two dots. The maximum distance between two points within a unit square is square root of 2.}
\end{22enumerate}

\subsection*{Checkpoint 2 - Call over a TA}

\newpage
\section*{Probability}
% \subsection*{Motivation and Definitions}

A \textbf{finite discrete probability space} is a pair $(S,\Pr)$ for some finite set $S$ and some function $\Pr:S\rightarrow \mathbb{R}$, where, for all $x \in S$, $\Pr(x)>0$, and $\sum_{x\in S}\Pr(x)=1$. The set $S$ is called the \textit{sample space}, and the elements of $S$ are called \emph{outcomes}. The function $\Pr$ is called the \textit{probability distribution function}. 

Where do these defintions come from? It's all about predictions. The sample space $S$ is supposed to represent a set of atomic outcomes---that is, the set of all things that can happen. All these outcomes are mutually exclusive.

An \textit{event} is a subset of the sample space $S$, containing some atomic outcomes. For example, if I flip a coin 3 times, the outcomes are HHH, TTT, \dots, so the sample space is $S=\{ \text{HHH}, \text{TTT}, \text{TTH}, \dots \}$, with $2^3=8$ outcomes. An event can be: \textit{We see exactly 1 H}, which we view as a set of outcomes $\{\text{HTT}, \text{THT}, \text{TTH}\}$. An event can also be the empty set: for example, the event \textit{We see 4 Hs} is the empty set, since it is not a possible outcome (it does not correspond to anything in the sample space). Similarly, the event \textit{We see an elephant} is the empty set for this probability space.

How many different events can the above sample space have? Since each event is a subset of the sample space, the number of events is equal to the number of possible subsets, which here is $2^8$ (since the sample space has 8 elements).

To understand the probability distribution function, let's think about our sample space as a box of area 1. Then, each outcome occupies a certain amount of \textit{area} within the sample space. That is, the area is \textit{distributed} among the outcomes according to the probability distribution function! The larger the area, the more likely the outcome is.

\includegraphics[scale=.7]{probability.png}

In this example, we have that each outcome occupies the same amount of area in the box---each is equally likely. Because we have 2 outcomes, we have that the probability of each outcome is $\frac{1}{2}$. 

Let's take a look at another way we could divide up the area of the box among the outcomes. 

\includegraphics[scale=.7]{probability2.png}

Here, the probability of it raining is $\frac{3}{4}$, so this outcome takes up $\frac{3}{4}$ of the area within the box. The remaining amount of area is $\frac{1}{4}$, which becomes the probability that it doesn't rain. 

\textbf{Task 5}
\begin{22enumerate}
  
\item If we have $n$ outcomes in $S$, and $\Pr$ assigns each outcome an equal amount of area within our box, what is the probability of a particular outcome? 

\note{Such a division of the space is called a \textit{uniform distribution}!}

\solu{1cm}{$1/n$}

\item Consider $S = \{\text{It's sunny}, \text{It's raining}\}$, where the probability distribution is 0.8 and 0.2, respectively. Under what constraints on the world is this sample space reasonable?

\solu{2cm}{If the world is only sunny or raining, then the sample space is fine. But, in the real world, we can have cloudy weather and snowy weather, too. A sample space should be \textit{exhaustive} and also \textit{mutually exclusive}, that is, exactly one of the outcomes has to happen each time we do the experiment.}


\item Is the following probability distribution valid? Experiment: flipping a biased coin. $S=\{H,T\}$, $\Pr(H)=1/3$, $\Pr(T)=1/3$.
\solu{2cm}{No, because the probability of all outcomes must add up to 1. }

\end{22enumerate}

\newpage
\subsection*{The Relationship Between Counting and Probability}

Questions about probability are often closely tied to questions about counting. Often, we need to count two things---the set of outcomes in our event, and the set of total possible outcomes. In a uniform distribution, the probability of getting an event $E$ is simply $\Pr(E) = |E|/|S|$.

\textbf{Task 6}
\begin{22enumerate}
  
\item Suppose $S$ is the set of all binary strings of length $n$, and $\Pr$ is a uniform distribution. What is the probability of getting a string with $k$ ones (where $0 \leq k \leq n)$? 

\solu{2cm}{There are $\binom{n}{k}$ strings with $k$ ones, and $2^n$ total strings, so the answer is $\binom{n}{k}/2^n$.}

\item For what value(s) of $k$ is the probability of getting a string with $k$ ones the largest? Try it out for some values and see if you find a pattern.

\solu{2cm}{When is $\binom{n}{k}$ largest? When $k$ is closest to $\frac{n}{2}$}

\item Explain how we can use $n$-ary (not just binary) strings to help us with other problems that involve, say, flipping a fair coin $m$ times, or rolling a die $m$ times. 

\solu{2cm}{We can form bijections to count things! For example, let 0 denote tail, head denote 1. Then, there is a bijection between the possible binary strings and the possible outcomes of $m$ coin flips. If we were using dice instead, though, we can use a senary (6 valued) string. }

\end{22enumerate}

\subsection*{Checkoff — call over a TA!}

\end{document}